\documentclass[12pt]{article}
\usepackage[utf8]{inputenc}
\usepackage[T1]{fontenc}
\usepackage[margin=1in]{geometry}
\usepackage{setspace}
\usepackage{xcolor}
\usepackage[colorlinks=true, urlcolor=blue!60!black, linkcolor=blue!60!black]{hyperref}
\setlength{\parindent}{0em}
\setlength{\parskip}{0.6em}
\renewcommand{\baselinestretch}{1.0}
\begin{document}
\begin{center}
{\LARGE \textit{San Antonio Independent School District v.\ Rodriguez}, 411 U.S.\ 1 (1973)\par}
\vspace{0.5em}
{\large Case Brief --- Two-Page Edition\par}
\end{center}

\section{Memory Jogger}
Unequal school funding from local property taxes does not violate equal protection: district wealth is not a suspect class, and education is not a fundamental right (\href{https://supreme.justia.com/cases/federal/us/411/1/}{Justia}).

\section{Detailed Case Facts}
Texas combined a state Foundation Program guaranteeing basic education with capped local property taxes. Property-poor Edgewood, about 90\% Mexican-American, spent \$356 per pupil at a high tax rate; affluent Alamo Heights spent \$594 at a lower rate. Edgewood parents led by Demetrio Rodriguez sued.

\section{Procedural History}
A three-judge district court held the system unconstitutional under strict scrutiny and rational basis review. Texas appealed directly.

\section{Judicial Votes}
Reversed, 5--4. Powell wrote for Burger, Stewart, Blackmun, and Rehnquist; Stewart concurred; Brennan, White (with Douglas and Brennan), and Marshall (with Douglas) dissented.

\section{Holding}
Strict scrutiny applies only to suspect classes and rights explicitly or implicitly guaranteed by the Constitution. Neither applies, and the system rationally serves local control.

\section{Analysis of Opinions}
\textbf{Majority:} No definable class of indigents suffered an absolute deprivation; Texas guaranteed free schooling. Education's importance does not make it fundamental, and the Court lacks expertise in school finance. \textbf{Stewart:} the system is ``chaotic and unjust'' but constitutional. \textbf{White:} capped tax rates gave poor districts no real local choice, failing rational basis review. \textbf{Marshall:} the Court should apply a ``spectrum of standards'' and calls local control a ``mere sham.''

\section{Examples: Future Applications}
\textbf{Same-Side:} school-bus fees without an indigency waiver (\textit{Kadrmas}). \textbf{Opposite-Side:} excluding undocumented children from school entirely (\textit{Plyler}). \textbf{Fence-Sitter:} a claim that a collapsing district denies basic literacy.

\section{Critique}
Judge Sutton credits the decision with spurring state-court reform; progressives say it entrenched ``separate and unequal'' schools.

\section{Key Quotations}
\begin{quote}
``The system of alleged discrimination and the class it defines have none of the traditional indicia of suspectness: the class is not saddled with such disabilities, or subjected to such a history of purposeful unequal treatment, or relegated to such a position of political powerlessness as to command extraordinary protection from the majoritarian political process.''
\end{quote}
\noindent --- Justice Powell, for the Court.

\end{document}
